\section{評価基盤の設計}\label{sec:methods}

\subsection{元データの点検と作業用データ}

入力は数値特徴量と分類ラベルを持つCSVである。まずヘッダー、欠損・非有限値、重複、同じ特徴量で正解ラベルが異なる組、クラス別件数を集計する。利用者は1行の意味、何を予測するか、予測時点、管理用IDや観測時刻、ラベルの由来を記録できる。これらは利用者が与える評価条件である。特徴量が実際に予測時点で使えるかは、記録の由来と照合する。

図\ref{fig:workflow}に、点検、評価、利用者別の出力をつなぐ運用を示す。たとえば10万行以上または20列以上を中継ファイルの作成を検討する目安とし、元CSVの読み直しや整形を繰り返す場合に作業用ファイルを再利用する。利用者の希望があれば、規模によらず圧縮ParquetまたはDuckDBファイルへ変換できる。現行版は利用者が変換を指定し、BP向け画面では圧縮Parquetを選べる。学習には選んだ列をCSVに書き出して渡す\cite{audit}。

学習に使う列を選び、行数、欠損、ラベル別件数、同じ入力でラベルが異なる組を再点検する。元CSVを保管し、評価対象の行は原則として全件を使う。欠損のある行を除外する場合は、基準と除外前後の件数を記録する。メモリや時間は学習時の行数、列数、モデルの設定から見積もる。予算を超える場合はモデル候補や計算設定を見直し、収まらなければ実行を保留する。

評価値と実行条件を保存し、同じ結果を用途に応じて提示する。データサイエンティスト（DS）向けにはコンソールの要約と予測・分割などの再分析用データを用意する。ビジネス利用者（BP）向けにはExcelで評価の対象件数、誤りの内訳、比較結果を示し、必要に応じて図を添える。

\begin{figure}[!t]
\centering
\begingroup
\setlength{\unitlength}{1mm}
\setlength{\fboxsep}{0pt}
\setlength{\fboxrule}{0.4pt}
\colorlet{FlowBorder}{Accent!65!white}
\colorlet{FlowStep}{Accent!5!white}
\colorlet{FlowChoice}{orange!12!white}
\colorlet{FlowReview}{black!6!white}
\colorlet{FlowOutput}{Accent!12!white}
% Arguments give the centre, width, height, fill and text, in that order.
\newcommand{\workflownode}[6]{%
  \put(#1,#2){\makebox(0,0){%
    \fcolorbox{FlowBorder}{#5}{%
      \parbox[c][#4mm][c]{#3mm}{\centering\fontsize{8.5pt}{11pt}\selectfont #6}}}}}
\newcommand{\workflowarrowhead}{%
  \begin{picture}(0,0)
    \put(-0.48,0.6){\makebox(0,0){\rotatebox{-51}{\rule{1.55mm}{0.4pt}}}}
    \put(0.48,0.6){\makebox(0,0){\rotatebox{51}{\rule{1.55mm}{0.4pt}}}}
  \end{picture}}
\begin{picture}(160,122)
  \color{FlowBorder}
  \linethickness{0.4pt}
  % Input and the choice of working format.
  \put(116,115){\line(1,0){8}}
  \put(124,115){\rotatebox[origin=c]{90}{\workflowarrowhead}}
  \put(80,110){\line(0,-1){5}}
  \put(80,105){\workflowarrowhead}
  \put(80,93){\line(0,-1){4}}
  \put(28,89){\line(1,0){104}}
  \put(28,89){\line(0,-1){4}}
  \put(80,89){\line(0,-1){4}}
  \put(132,89){\line(0,-1){4}}
  \put(28,85){\workflowarrowhead}
  \put(80,85){\workflowarrowhead}
  \put(132,85){\workflowarrowhead}
  % The three formats converge on the same preparation step.
  \put(28,75){\line(0,-1){4}}
  \put(80,75){\line(0,-1){4}}
  \put(132,75){\line(0,-1){4}}
  \put(28,71){\line(1,0){104}}
  \put(80,71){\line(0,-1){2}}
  \put(80,69){\workflowarrowhead}
  \put(80,59){\line(0,-1){5}}
  \put(80,54){\workflowarrowhead}
  \put(113,49){\line(1,0){5}}
  \put(118,49){\rotatebox[origin=c]{90}{\workflowarrowhead}}
  \put(80,44){\line(0,-1){4}}
  \put(80,40){\workflowarrowhead}
  % Both outputs use the saved evaluation.
  \put(80,28){\line(0,-1){5}}
  \put(40,23){\line(1,0){80}}
  \put(40,23){\line(0,-1){4}}
  \put(120,23){\line(0,-1){4}}
  \put(40,19){\workflowarrowhead}
  \put(120,19){\workflowarrowhead}

  \color{black}
  \workflownode{80}{115}{72}{10}{FlowStep}{\textbf{元CSVを保管して点検}}
  \workflownode{142}{115}{35}{10}{FlowReview}{不備があれば\\点検結果を返す}
  \workflownode{80}{99}{72}{12}{FlowChoice}{\textbf{中継データを作るか}\\規模・再利用・利用者指定}
  \workflownode{28}{80}{45}{10}{FlowStep}{元CSVを使う}
  \workflownode{80}{80}{45}{10}{FlowStep}{圧縮Parquetに変換}
  \workflownode{132}{80}{45}{10}{FlowStep}{DuckDBファイルに変換}
  \workflownode{80}{64}{88}{10}{FlowStep}{学習用の列を選び、欠損と件数を再点検}
  \workflownode{80}{49}{66}{10}{FlowChoice}{\textbf{対象行を保ったまま学習できるか}}
  \workflownode{138}{49}{39}{12}{FlowReview}{設定を見直す\\収まらなければ保留}
  \workflownode{80}{34}{104}{12}{FlowStep}{\textbf{保留データを確保し、学習側の交差検証で順位付け}\\保留評価と不確かさの結果を保存}
  \workflownode{40}{11}{72}{16}{FlowOutput}{\textbf{DS向け}\\コンソール要約・再分析用データ}
  \workflownode{120}{11}{72}{16}{FlowOutput}{\textbf{BP向け}\\Excel要約・必要に応じて図}
\end{picture}
\endgroup
\caption{データ点検からDS・BP向け出力までの運用フロー。中継形式と学習条件は利用者が指定する。}
\label{fig:workflow}
\end{figure}


\begin{table}[htbp]
\centering
\small
\caption{工程と、あとから確認する主な記録}
\label{tab:pipeline}
\begin{tabularx}{\linewidth}{@{}p{31mm}YY@{}}
\toprule
工程 & 確認する内容 & 保存する情報 \\
\midrule
元CSVの点検 & 欠損、重複、ラベル別件数と評価条件 & 入力のハッシュ、点検Excel、JSON \\
抽出と再点検 & 使う特徴量、行数、クラス構成 & 抽出条件、抽出CSV、再点検の結果 \\
モデル選定 & 学習側の共通分割と候補ごとの指標 & 設定、検証行、平均と区間ごとの成績 \\
保留テスト & 選定モデルと単純ルールの結果 & 予測、正解、誤りの件数、評価区間の適用範囲 \\
\bottomrule
\end{tabularx}
\end{table}

\subsection{学習側の候補比較と単純ルール}

教師あり分類の基本候補は、決定木、ランダムフォレスト（100木）、ロジスティック回帰（最大1,000反復）、SVM、5近傍法である。各検証のスケーラーは学習部分だけから推定し、全候補に同じ分割を与える。既定の選定指標はクラスごとのF1の算術平均である。クラス集合を \(\mathcal{C}\) とし、クラス \(c\) の真陽性・偽陽性・偽陰性件数を \(\mathrm{TP}_c,\mathrm{FP}_c,\mathrm{FN}_c\) とすると、
\begin{equation}
\macroF = \frac{1}{|\mathcal{C}|}\sum_{c\in\mathcal{C}}
  \frac{2\,\mathrm{TP}_c}{2\,\mathrm{TP}_c+\mathrm{FP}_c+\mathrm{FN}_c}
\label{eq:macro}
\end{equation}
である。本実装では分母が0の項を0とし、テストや月別集計でクラス集合を固定する。ラベル別の検出率は該当する正解行が0件なら空欄とする。

5回の検証の平均で候補を選び、同率判定の許容値は絶対差 \(10^{-12}\) とする。同率の場合は宣言した候補順を使う。選定に使った検証スコアには選択の影響がある\cite{varma}。このため、学習側の全範囲で再学習した一つのモデルを別に取った保留テストへ適用する。単純ルールの閾値や多数派は学習側で求める。比較では、精度指標に加え、正解件数、混同行列またはラベル別の見逃し・誤検知を確認する。繰り返し分割で順位を調べる試行は、固定の5回で代表モデルを選ぶ経路と区別して保存する\cite{ranking}。

\subsection{行・グループ・日付を単位とする分割}

\begin{description}
  \item[行単位] クラス構成を保つ層化分割で約20\%の行を保留し、学習側を共通の層化5分割で評価する。
  \item[グループ単位] 指定したID、または選んだ数値特徴量の完全一致をグループとする。学習と保留、および学習側の各検証で、同じグループを左右に分けない。層化グループ分割を使い、保留と学習側の各検証区間に全クラスが現れる分割を選ぶ\cite{groupkfold,audit}。管理用IDは特徴量から除く。数値入力の一致は観測値の定義であり、同一人物を識別する情報とは限らない。
  \item[日付順] 観測日のうち新しい日から約20\%を取り置く。学習側を観測日順の6区間に分け、最初の区間で学習し、次の区間で評価する。その区間を学習範囲に加えて同じ操作を5回行う。同じ観測日の行はまとめて扱い、日付列は特徴量から除く。時差のある時刻はUTCの日付でそろえ、時差指定のない時刻は入力の暦日を使う。各期間に必要なクラスがそろわない場合は、設定した評価を実行できないことを表示する\cite{audit}。
\end{description}

日付順で学習を進めるには、学習範囲に入れた正解ラベルが次の評価期間の開始時点までに得られることが必要である。今回の方式では観測日で分割し、ラベルの確定が遅れる場合の待機期間は設計課題に残した。日付順と顧客を分ける条件の両方を課す評価も後続で扱う。

\subsection{固定モデルの区間、計算精度、偏りの点検}

標本の不確かさには、選定・学習済みのモデルを固定して保留行の正解と予測を対にしたブートストラップを使う\cite{efron1979}。行単位の層化方式では、正解クラスごとに観測件数を保って復元抽出する。グループ分割では、観測したグループ数と同数のグループを復元抽出し、選ばれた各グループの行をまとめて用いる。行に基づく分類指標では、抽出されたグループのサイズに応じて寄与する行数が変わる。

\(B\) 回の再標本化で指標 \(T^{*(1)},\ldots,T^{*(B)}\) を求める。標準誤差の集計は反復値の標本標準偏差、95\%区間は経験分布の2.5と97.5百分位点とする。基本比較と本稿の行・グループの主要例では \(B=1{,}000\) を用いた。固定モデルに条件付けた値なので、収集、特徴量選択、モデル選定、再学習を繰り返した場合の変動は別の評価を要する。行やグループの代表性も、標本設計と実データの由来から検討する。

同じテスト対象で二つの方法を比較する追加の診断では、同じ再標本化の添字を双方に使い、差 \(T_A^*-T_B^*\) を集計する。本稿の日付順評価では性能値と月別・ラベル別の内訳を報告する。連続する観測日の依存を保つ区間は未算出である。

計算精度の確認には、同じ混同行列から有理数で求めた指標と浮動小数点での算出値の照合や、入力をFloat32へ丸めた場合の予測比較を用いる。入力の校正ずれを仮定した感度分析は、指定したずれに対する応答を示す。実際の系統的な偏りの有無や大きさには、原票、独立した測定、ラベル再判定などの情報が要る\cite{ranking}。
